\documentclass[11pt,letterpaper]{article}

\usepackage[T1]{fontenc}
\usepackage[spanish,es-nodecimaldot]{babel}
\usepackage[margin=2.5cm]{geometry}
\usepackage{amsmath,amssymb}
\usepackage{enumitem}
\usepackage{graphicx}
\usepackage[hidelinks]{hyperref}
\usepackage{multicol}

\title{
    Investigación N\textsuperscript{o} 2 RM3100 \\
    \large Información sobre distanciamiento entre el sensor y la electrónica
}

\author{Yerko Jelcic Iturra}
\date{\today}

\begin{document}

\maketitle

\begin{abstract}
La pregunta que se busca resolver en este informe es sobre la necesidad de separar el cabezal del sensor RM3100 de los componentes electrónicos.
Ya que, al investigar, se encontró que hay dos variables que se deben considerar al momento de diseñar el sistema integrado, siendo el primero la distancia y el segundo la minimización del \textit{current-loop area} (área del bucle de corriente). Algunos proyectos comparables, como el de la \textit{HamSCI's Personal Space Weather Station}, que divide el diseño en una ``placa remota'' que contiene solo el RM3100, colocada en el exterior de un tubo de PVC y conectada a un Raspberry Pi en el interior mediante un par trenzado blindado que transmite $I^2C$, no SPI directo.\\
Por otro lado, un nodo compacto totalmente integrado puede llegar a ser posible y esta validado en aplicaciones espaciales \textit{(CubeSat "Quad-Mag" bus-mounted boards)}, pero solo cuando se combina con algoritmos de cancelación de ruido y/o un capo de plataforma estable y calibrado. Además, las configuracioens integradas de diversos proyectos con el sensor RM3100, junto con un ESP32/Wi-Fi y fuentes de alimentación conmutadas, suelen presentar ruido de entre 50 y 200 [nT], que corresponde a aproximadamente 10 veces peor que la especificación del sensor en $\sim$ 15 [nT].
\end{abstract}

\begin{multicols}{2}

\section{Key Findings}
Dentro de los proyectos investigados, se tieen que lo más razonable es separar los sensores por diseño. El magnetómetro 


\end{multicols}
\end{document}